\LevelZero{Q\&A}\label{ch:q&a}


\LevelOne{Where is the configuration files of fuji?}
As a convention, all the files are placed in \ttt{config/fuji/} directory.


\LevelOne{What is .json file?}
A json file is a text file, whose file name normally ends with \tbf{.json}.


\LevelOne{How can I edit a configuration file?}
To ensure the \ttt{readibility} and \ttt{compatibility}, most of the files are saved as \ttt{pure text format}.
You can open them with a \ttt{text editor}.
Remember to issue \cmd{/fuji reload} command to apply your modification.

\begin{tips}{Use a moderen text editor.}
    Some files may have a large number of lines, so it's highly recommended to use a \ttt{modern} text editor, which can highlight symbols and reveal the structure of the file, such as:
    \begin{enumerate}
        \item \href{https://code.visualstudio.com/}{Visual Studio Code}
        \item \href{https://vscode.dev/}{Visual Studio Code - Web Online Editor}
        \item \href{https://www.vim.org/}{Vim}
        \item \href{https://www.gnu.org/software/emacs/}{Emacs}
        \item \href{https://www.sublimetext.com/}{Sublime Text}
    \end{enumerate}
\end{tips}


\LevelOne{What is .dat file?}
The file whose name ends with \tbf{.dat} are the \ttt{vanilla minecraft NBT format file}.
To open such a file, you need to use a \ttt{NBT Editor}, such as \ttt{NBTExplorer}.


\LevelOne{How to update fuji to a new version?}

Please follow the steps:
\begin{description}
    \item[1. Backup the data] {Back up the \ttt{config/fuji} directory.}
    \item[2. Test the new version in your test-environment] {Put the new version of fuji into \ttt{mods/} directory, start the server, and adjust the configuration to your desired effect.}
    \item[3. Apply the changes to production-environment] {Now, it's ready to apply the changes to your production-environment.}
\end{description}

\begin{warn}{Don't test new changes directly in your production-environment}
    It's highly recommended to setup a test-environment for a network-maintainer, so that you can test and tweak installed mods into the desired effect, and avoid un-expected situations.
\end{warn}


\LevelOne{Fuji conflits with one of my mods.}
You can \ttt{disable} that \ttt{conflicting module} and re-start your server.
If possible, create an issue at \issueurl, so that we can solve it later.

\LevelOne{Can I ask the forge or neoforge support?}
We have no plan for forge or neoforge platform.
However, you can try running this mod via \tbf{sinytra-connector} mod.
It's tested in the following environment (105/105 modules works):
\begin{shcode}
    Loader: Forge mod loading, version 47.4.0, for MC 1.20.1 with MCP 20230612.114412
    Mods:
    - connector-1.0.0-beta.46+1.20.1
    - forgified-fabric-api-0.92.2+1.11.12+1.20.1
    - fuji-fabric-7.4.0-b93ba03936-mc1.20.1
\end{shcode}

If you encounter issues in forge or neoforge environment, you can still open an issue in github.
We provide the support for forge and neoforge, if it's possible and handy.

\LevelOne{How can I report bugs or suggest new features?}
You can create an issue at \issueurl



